% Generated draft: requires mathematical and source review.
\documentclass[11pt]{article}
\usepackage[margin=25mm]{geometry}
\usepackage{fontspec}
\setmainfont{Latin Modern Roman}
\usepackage{amsmath,amssymb,mathtools,unicode-math}
\setmathfont{Latin Modern Math}
\usepackage{enumitem,needspace,xurl,hyperref}
\setlength{\parindent}{0pt}
\setlength{\parskip}{5pt}
\setlength{\emergencystretch}{3em}
\newcommand{\N}{\mathbb{N}}
\newcommand{\Z}{\mathbb{Z}}
\newcommand{\Q}{\mathbb{Q}}
\newcommand{\R}{\mathbb{R}}
\newcommand{\C}{\mathbb{C}}
\newcommand{\F}{\mathbb{F}}
\newcommand{\A}{\mathbb{A}}
\newcommand{\PP}{\mathbb{P}}
\newcommand{\E}{\mathbb{E}}
\newcommand{\eps}{\varepsilon}
\newcommand{\dd}{\,\mathrm{d}}
\begin{document}
\section*{UnsolvedMath ID: 6200001}
\textbf{Source problem code:} catalogue entry
\subsection{Definitions and mathematical statement}
A finitely generated group \(G\) is word-hyperbolic if one, equivalently any, Cayley graph for a finite generating set has uniformly thin geodesic triangles. Its Gromov boundary \(\partial G\) consists of geodesic rays modulo finite Hausdorff distance, with the standard boundary topology. The first question asks for a characterization of the topological spaces homeomorphic to \(\partial G\) as \(G\) ranges over word-hyperbolic groups. No torsion-free or one-ended restriction is stated.

The specified subproblem asks which dimensions occur for ``stable Menger spaces''. The supplied research summary interprets these as universal Menger compacta. For that interpretation, let \(\mu^k\) be the standard \(k\)-dimensional Menger compactum: start with \([0,1]^{2k+1}\), subdivide each coordinate into thirds, retain subcubes for which at most \(k\) coordinates are middle thirds, and repeat the same operation in every retained cube; \(\mu^k\) is the intersection of the retained compact sets. Determine
\[
 \{k\in\mathbb Z_{\ge0}:\exists\text{ a word-hyperbolic group }G
 \text{ with }\partial G\cong\mu^k\}.
\]
Here \(\cong\) denotes homeomorphism, not a specified visual-metric isometry or quasisymmetric equivalence. The word ``stable'' is not defined in the excerpt, so identifying it with this standard universal compactum requires source confirmation.

The supplied summary reports realization for \(k=1\), the Menger curve, and \(k=2\), and says higher-dimensional realization and the full boundary classification remain open. No realizing groups or original proof references are supplied. These are dated progress attributions; the 2005 source's problem label and the current registry assessment are not an independently verified literature review.
\subsection{Short English statement}
Which topological spaces are boundaries of word-hyperbolic groups, and in particular which universal Menger compacta occur? The record reports dimensions one and two realized, while the precise meaning of its term ``stable Menger space'' and the cited progress need source review.
\subsection{Sources}
\begin{itemize}
\item Local registry: \url{data/unsolvedmath/problems.json}, permanent UnsolvedMath ID \(6200001\), source code catalogue entry; supplied statement, background, and research summary, updated 2026-08-06.
\item Misha Kapovich, \emph{Problems on Boundaries of Groups and Kleinian Groups} (2005), Problem 1, PDF page 2: \url{https://www.math.ucdavis.edu/~kapovich/EPR/problems.pdf}. Supplied normalized PDF-text extraction accessed 2026-07-29; current status marked as needing review.
\item The summary reports Menger-boundary realizations in dimensions \(1\) and \(2\), but supplies no original realization papers or examples.
\end{itemize}
\end{document}
